% ID: SEL-D0064-OA-E02
% ESTADO: validada. S064.
\textbf{a)} $N(t)=4-(t-2)^2$. El cuadrado se anula en $t=2$, de modo que el máximo es de $4$ pacientes de media a las $\boxed{19{:}00}$.

\textbf{b)} $N(t)=t(4-t)$ se anula al abrir y cuatro horas después. Entre ambos instantes es positiva. El consultorio cierra a las $\boxed{21{:}00}$.

\textbf{c)} La gráfica con $N(t)\ge0$ es el arco de parábola entre $(0,0)$ y $(4,0)$, con vértice $(2,4)$.
\begin{center}
\setlength{\unitlength}{0.9mm}
\begin{picture}(94,55)(-9,-7)
\qbezier(0,0)(41,0)(82,0)
\qbezier(0,0)(0,22)(0,44)
\qbezier(0,0)(36,64)(72,0)
\put(83,0){$t$}\put(2,44){$N(t)$}
\put(-6,-5){\scriptsize $0$}
\put(33,35){\scriptsize $(2,4)$}
\put(70,-5){\scriptsize $4$}
\end{picture}
\end{center}
